\documentclass[11pt,a4paper,english,oneside]{book}
\usepackage{etex} %Because of many packages --> Extended TeX.
\usepackage[left=1in, right=1in]{geometry} %Helps to structure the paper layout.
\usepackage[Lenny]{fncychap} %Design of the thesis.
\usepackage[utf8]{inputenc} %Due to vowels.
\usepackage{babel} %Define the language style.
\usepackage{dsfont} %Nice style for the indicator function.
\usepackage{fancyhdr} %To customize the headers and footers.
\usepackage{booktabs} %In case you need \cmidrule or \addlinespace in tables.
\usepackage[hang,bottom,stable,multiple]{footmisc} %Style of footnotes.
\usepackage{appendix} %For the \appendixpage command.
%Load some mathematical packages.
\usepackage{amsmath}
\usepackage{amsfonts}
\usepackage{amsmath}
\usepackage{amssymb}
\usepackage{mathtools}
%\usepackage[sort,round]{natbib} %For the bibliography.
\usepackage[natbibapa]{apacite} %for citing and bibliography, by Corinne
\usepackage{etoolbox} %To remove the page number on \appendixpage.
\usepackage{amsthm} %For theorems, definitions etc.
\usepackage{thmtools} %For theorems, definitions etc.
\usepackage{setspace} %Use double spacing.
\usepackage{lipsum} %For the \lipsum command to generate a text.
\usepackage{datetime} %For the specification of the date.
%\usepackage{tocloft} %The ToC, LoF and LoT each appear not necessarily on a new page.
\usepackage{graphicx,listings,xcolor,textcomp} %For the graphics, listings etc.
\usepackage{mcode} %To implement a Matlab code.
\usepackage[margin=10pt, font=small, labelfont=bf, labelsep=endash]{caption} %Customize the captions.
\usepackage{chngcntr} %To use counterwithout.
\usepackage{epstopdf} %For inserting .eps files into the document.
\usepackage{hyperref} %Must be loaded at the end.
\usepackage{xparse} %Load for \NewDocumentCommand command.
\usepackage{cleveref} %For the command \cref, load after hyperref.
\usepackage{arydshln} %Due to the capability to draw horizontal/vertical dash-lines.
\usepackage{array,hhline} %To create tables and matrices.
\usepackage{rotating} %To rotate a table.
\usepackage{tabularx} %An extended version of tabular.

%Setup of the reference links.
\hypersetup{
     colorlinks=false,
     linkcolor=blue,
     citecolor=blue,
     filecolor=magenta,
     urlcolor=blue}

%Define some reasonable margins.
%\setlength{\textwidth}{6.6in}
\setlength{\textheight}{9.3in}
\setlength{\topmargin}{-0.1in}
%\setlength{\oddsidemargin}{0in}
\setlength{\parskip}{1mm}

\bibliographystyle{apacite} %Reference style.
\allowdisplaybreaks[1] %Page breaks of equations are allowed, but avoided if possible. 2-4 more relaxed.

%New command for the logo.
\newcommand*{\plogo}{\includegraphics[scale=0.7]{HSG_logo}}

%New command for the differential d to have an ordinary d.
\makeatletter
  \newcommand{\ud}{\mathrm{d}}
\makeatother

%Remove page number on \appendixpage. Use the package 'etoolbox'.
\makeatletter
\patchcmd{\@chap@pppage}{\thispagestyle{plain}}{\thispagestyle{empty}}{}{}
\makeatother

%Declare Definitions, Theorems etc.
%%%%%%%%%%%%%%%%%%%%%%%%%%%%%%%%%%%%%%%%%%%%%%%%%%%%%%%%%%%%%%%%%%%%%%%%%%%%%%%%%%%%%%%%%%%%%%%%%%%%%%%%%%%%%%%%%%%
\declaretheorem[style=definition,qed=$\blacktriangleleft$, numberwithin=chapter]{remark} %additional options; numberwithin=,..., see 'Thmtools' Users’ Guide
\declaretheorem[style=definition,qed=$\triangle$,numberwithin=chapter]{definition}
\newtheorem{ass}{Assumption}[chapter]
\newtheorem{prop}{Proposition}[chapter]
\newtheorem{lemma}{Lemma}[chapter]
\declaretheorem[style=definition,qed=$\perp$,numberwithin=chapter]{example}
\newtheorem{theorem}{Theorem}[chapter]
\newtheorem{coroll}{Corollary}[chapter]
%%%%%%%%%%%%%%%%%%%%%%%%%%%%%%%%%%%%%%%%%%%%%%%%%%%%%%%%%%%%%%%%%%%%%%%%%%%%%%%%%%%%%%%%%%%%%%%%%%%%%%%%%%%%%%%%%%%

%\clubpenalty = 10000 %von Corinne
%\widowpenalty = 10000
%\displaywidowpenalty = 10000


%Readjust the numbering.
%\counterwithout{footnote}{chapter}
\numberwithin{equation}{chapter}

%\setlength{\parindent}{0cm} %Uncomment this if you don't want to have indents.

%----------------------------------------------------------------------------------------
%	TITLE PAGE
%----------------------------------------------------------------------------------------
\newcommand*{\titleGP}{\begingroup %Create the command for including the title page in the document.
\centering %Center all text.
\vspace*{\baselineskip} %White space at the top of the page.
\plogo\\[2\baselineskip] %University Logo.
\rule{\textwidth}{1.6pt}\vspace*{-\baselineskip}\vspace*{2pt} %Thick horizontal line.
\rule{\textwidth}{0.4pt}\\[\baselineskip] %Thin horizontal line.
{\LARGE Narratives in Finance }\\[0.2\baselineskip] %Title.
\rule{\textwidth}{0.4pt}\vspace*{-\baselineskip}\vspace{3.2pt} %Thin horizontal line.
\rule{\textwidth}{1.6pt}\\[2\baselineskip] %Thick horizontal line.
\scshape %Small caps.
Master's Thesis\\[2\baselineskip]
Submitted in partial fulfillment of the requirements for the degree of Master of Arts in Quantiative Economics and Finance \par
\vspace*{2\baselineskip}
Author\\
{\Large Corinne Knöpfel \\ [5pt]
 }
Wiesen 2488, Herisau \\[5pt]
11-613-676\\[5pt]
corinne.knoepfel@student.unisg.ch \\


\vspace*{2\baselineskip}
Supervisor\\
{\Large Dr. Peter Gruber\\[5pt]%\small Hans Vontobel Professor of Financial Engineering\\[5pt]
\small Institute of Finance\\[5pt]Universit\`{a} della Svizzera Italiana\par}
%\vspace*{2\baselineskip}
%Assistant\\
%{\Large [ Name ] \par}
\vfill
{\scshape Date of Submission: \today } \\[0.3\baselineskip]
\endgroup}

%Special header and footer style for the Abstract and other special pages.
\fancypagestyle{firststyle}{%
  \fancyhf{}%
  \renewcommand{\headrulewidth}{0pt}
  \fancyfoot[C]{\thepage}
}


%Customize headers and footers.
\pagestyle{fancy}
\fancyhead[R]{\thepage}
\fancyhead[L]{\rightmark} %section in header
%\fancyhead[L]{\leftmark}  %chapter in header
%\fancyfoot[L]{[ Your Name ]}
\fancyfoot[C]{}
%\fancyfoot[R]{[ Running Master Thesis Title ]}

%Define the signature line with dots.
\NewDocumentCommand \dotbox {o O{.5\linewidth} m O{3ex} O{\linewidth}}
{
  \begin{minipage}{7cm}
    \makebox[7cm][l]{\,.\dotfill}
    \\
    \makebox[7cm][l]{\,#3}
  \end{minipage}
}

\begin{document}
\thispagestyle{empty}
\titleGP
\newpage
\doublespacing
\setcounter{page}{1}
\pagenumbering{Roman}
%\section*{Task Assignment}
\thispagestyle{firststyle}
%\pagestyle{firststyle} %from here on it uses pagestyle firststyle (for content, etc.)


%\newpage
\chapter*{Abstract}

{\pagestyle{firststyle}
\tableofcontents
\cleardoublepage
}

\listoffigures
\listoftables
%\setcounter{rememberpage}{\value{page}} %in case we want to go back to Roman numbering in the Appendix

\newpage


\pagenumbering{arabic}
% \part{[ Part title ]}
\chapter{Introduction}
 

\chapter{Monetary Policy and Interest Rates} \label{MonetaryPolicy}
 
\noindent However little understood, the relationship between monetary policy and market interest rates is undeniable. Interest rates of all maturities react to changes in monetary policy, creating opportunities and risks for traders, challenges for policy makers, and puzzling effects for academics to study (\citealt[p. 1594]{Ellingsen.2001}). 

Target rate changes in particular have an impact on the bond market and on interest rates (\citealt[p. 332]{Cook.1989}). %Overall, there is an invested interest in how monetary policy impacts yield curve movement.
Yet, the understanding of yield curve movements is incomplete at best. On average, the relationship between monetary policy and interest rates appears to be positive: An increase in the central bank's target rate leads to an increase in the interest rates of all maturities. However, there are many instances where this simple rule has proven false and interest rates of long maturities fell in response to an increase in the central bank's rate \citep[p. 1594]{Ellingsen.2001}. 

Chapter \ref{SensitivityPuzzle} gives an account of the puzzle posed by the inconsistent response of long-term rates, Chapter \ref{ExistingResearch} touches on previous research on the topic, and Chapter  \ref{NewInsights} outlines how an investigation of narratives might be able to shed light on this puzzle.




\section{Excess Sensitivity Puzzle} \label{SensitivityPuzzle}

\cite{Cook.1989} analyzed financial data from the late 70s and found that the U.S. Federal Reserve (Fed), by setting the target for the federal funds rate, had a strong influence on interest rate movements. While short-term rates reacted particularly strongly, changes in the target rate also caused small but significant movements in long-term rates. %According to \citeauthor{Cook.1989}, the positive relationship between the funds rate and market interest rates is evidence in favor of the expectation theory of the term structure of interest rates. They are convinced, that ***explain this and then contrast*** // 

 
It is not surprising that short-term rates follow the target rate closely, after all the Fed keeps the overnight rate close to the target and thus directly influences the one-month rate \citep[p. 1]{Ellingsen.2003}. The movements of the long-term rates are more ambiguous. \citet[p. 343--346]{Cook.1989} interpret the fact that, on average, 10-year and 30-year bonds co-move with the short-term rates as evidence for the expectation theory of the term structure of interest rates. According to the expectation theory, long-term rates are equal to short-term rates over the same period of time plus a term premium. Thus, an increase in the short-term rates is expected to drive up long-term rates as well, but to a lesser extent \citep[p. 1594]{Ellingsen.2001}.


%In theory, the effect of a target rate change could be offset by an adjustment of inflation expectations, pushing long-term rates in the opposite direction \citep{Ellingsen.2003}. Yet, \cite{Cook.1989} found clear evidence that, on average, even 10-year and 30-year bonds co-move with the target rate and the short-term rates. 

To \cite{Romer.2000}, on the other hand, the response of long-term interest rates presents a puzzle. They argue that standard theory predicts a drop in inflation as short-term rates rise, which ought to lead to a reduction in long-term rates. The opposite can be observed, however: Interest rates for all maturities typically rise following an increase in the target rate. \cite{Romer.2000} explain this anomaly with information-asymmetry between the Fed and the general public. They find evidence that the Fed is in possession of private information, which it reveals to other market participants through its monetary policy. In response, market participants adjust their inflation expectations upwards, causing long-term rates to rise.


%Yet, \cite{Cook.1989} found clear evidence that, on average, even 10-year and 30-year bonds co-move with the target rate and the short-term rates. 
%interpret their finding as supportive of the expectations theory of the term structure.2 The expectations theory says that a long interest rate should be equal to the average of short interest rates over the same period of time plus a term premium; thus an increase in the first couple of short rates should drive up the long rate too, but by less.


Dissecting the interest rate response in more detail paints an even complexer picture. While the yield curve shifts upwards on average, Skinner and Zettelmeyer (1995) found a number of occasions where an adjustment to the target rate caused the yield curve to tilt: Long and short rates responded by moving in opposite directions (as cited in \citealt[p. 1]{Ellingsen.2003}). Skinner and Zettelmeyer came to the conclusion that these were not singular occurrences, but that such tilts made up a considerable portion of the yield curve responses and could be observed in all four of the big economies they studied, that is in France, Germany, the United Kingdom, and the United States (as cited in \citealt[p. 1594]{Ellingsen.2001}). An example is the yield curve movement in 1994, where interest rates of long maturities fell after the Fed announced an increase in its target rate \citep[~p. 1594]{Ellingsen.2001}.

 

\citet[p. 425]{Gurkaynak.2005} argue that long term interest rates, unlike short term interest rates, should not react to unexpected macroeconomic events or monetary policy changes, since such changes have only a transitory effect on future interest rates. He says that can be explained by market expectations / opposition to Cook and Han and romer and romer who see this as aberation


s standard macroeconomic models predict that sh that short term interest rates return to their steady state 
we demonstrate that longterm forward rates move significantly in response to the unexpected components of many macroeconomic data releases and monetary policy announcements.

Their model can account for the observed negative response of long-term forward rates to monetary policy innovations by allowing for central bank private information about the inflation target. However, they are not able to explain the strong positive response of long-term yields to monetary policy found in most empirical studies nor the time-varying response of long rates to monetary policy.


%
This strong response of long-term interest rates to monetary policy innovations is typically difficult to replicate using standard macro models (see also G¨urkaynak et al., 2003).





Interest rate sensitivity is a measure of how much the price of a fixed-income asset will fluctuate as a result of changes in the interest rate environment. Securities that are more sensitive have greater price fluctuations than those with less sensitivity. This type of sensitivity must be taken into account when selecting a bond or other fixed-income instrument the investor may sell in the secondary market

Generally, the longer the maturity of the asset, the more sensitive the asset to changes in interest rates. Changes in interest rates are watched closely by bond and fixed-income traders, as the resulting price fluctuations affect the overall yield of the securities. Investors who understand the concept of duration can immunize their fixed-income portfolios to changes in short-term interest rates.


\section{Existing Research} \label{ExistingResearch}

E2001 is a narrative approach

\begin{quote}It has often been noted that the response of long-term interest rates to monetary policy innovations differs from occasion to occasion: sometimes long rates move in the same direction as the policy innovation, sometimes they move in the opposite direction. Most models of monetary policy cannot account for this puzzling behavior of long-term interest rates. In our previous work, we have shown that such a behavior is easily explained in a model where the central bank has private information about economic shocks and its own preferences or targets.~\cite{Ellingsen.2004}\end{quote}
%
\begin{quote}(2001) find that the yield curve response to monetary policy innovations depends crucially on the interpretation of bond market participants of the reasons behind the policy move.
	
	The intuition behind these results is straightforward. When supply or demand shocks cannot be directly observed, any unanticipated increase in the central bank’s policy rate is interpreted as a response to an unobserved inflationary shock. As the central bank is expected to counteract this inflationary impulse by tightening policy for some time, interest rates of all maturities increase as market participants update their expectations of the future path of the short rate. If, on the other hand, shocks are observable, but central bank preferences or objectives are not, an unanticipated tightening of policy is interpreted as a shift to a more inflation averse policy. Such a shift will imply a period of tighter policy than previously expected, but a quicker return to a neutral stance. Thus, short-term rates will increase in response to the policy innovation, while longer rates fall.~\cite{Ellingsen.2004}\end{quote}

\begin{quote}In Ellingsen and S¨oderstr¨om (2003) we test these theoretical predictions by classifying policy moves in the U.S. as endogenous or exogenous using reports in the \textit{Wall Street Journal. }The results are illustrated in Figure 3. Panel (a) reiterates the results from Figure 1, showing the estimated response of the yield curve to changes in the three-month T-bill rate (our measure of policy innovations) on all days when the Federal Reserve’s target for the federal funds rate was changed from October 1988 to December 2001.8~\cite{Ellingsen.2004}\end{quote}

\begin{quote}after policy moves classified as endogenous, interest rates of all maturities tend to move in the same direction, but after moves classified as exogenous, long and short rates move in opposite directions.10~\cite{Ellingsen.2004}\end{quote}

Idee: es hängt von der Interpretation ab, von der Narration die darum herum aufgebaut wird

\cite{Krosigk.2017} tried to replicate the results using text mining techniques but was unable to do so for the years xxx.

Ellingsen, Söderström and Masseng (2003) run two regressions to test their model of market interest rates and target adjustments. First, they test their theory whether the relationship between long and short rates differs on policy days and non-policy days. For that purpose, they estimate the regression \cite[p. 30]{Krosigk.2017}

$\Delta i^n_t = \alpha + (\beta_n^{NP}d_t^{NP} + \beta_n^{P}d_t^{P})\Delta i^{3m}_t + v_t^n$

H1: for large $n$: $\beta_n^{P}< \beta_n^{NP}$

Secondly, Ellingsen, Söderström and Masseng (ibid.) investigate whether the long and short rates behave differently on policy days classified as endogenous or exogenous as well as whether non-policy days have a similar impact as endogenous policy days. \cite[p. 30]{Krosigk.2017}

$\Delta i^n_t = \alpha + (\beta_n^{NP}d_t^{NP} + \beta_n^{Ex}d_t^{Ex} + \beta_n^{End}d_t^{End})\Delta i^{3m}_t + v_t^n$

H2: for large $n$: $\beta_n^{Ex}<0< \beta_n^{End}$

\section{New Insights Through Narrative Research} \label{NewInsights}

Narratives / interpretation plays a big role in how people react to financial facts.

\begin{quote}In joint work with Goetzmann and Kim (2016), using data from a questionnaire survey I have been conducting with institutional investors and high-income Americans since 1989, we found that these people generally have exaggerated assessments of the risk of a stock market crash, and that these assessments are influenced by the news stories, especially front page stories, that they read. One intriguing finding was that an event such as an earthquake could influence estimations of the likelihood of a stock market crash.~\cite[S.~974]{Shiller.2017}\end{quote}

also kann Textauswertung etwas dazu beitragen, es geht um das Verständnis zu Narrativen - aber jetzt die Frage der Zeitdimension: die Kurse sind innerhalb einer halben Stunden verändert - question: is the narrative really driving the change still or is this rather a case of already observing the result -> we can't explain y by knowing y!

So, use NLP techniques and then model like thus: aber jeder regressor erh;ht nat[rlich das R2]

$\Delta i^n_t = \alpha + (\beta_n^{NP}d_t^{NP} + \beta_n^{N1}d_t^{N1} + \beta_n^{N2}d_t^{N2} + + \beta_n^{N3}d_t^{N3} + ...)\Delta i^{3m}_t + v_t^n$


\cite{Cook.1989} use Wall Street Journal articles, even though they mention the speculative wording of the journals, they try their best ot isolate the facts in the articles, completely ignoring the manner in which the facts are presented> giving example of wording
Even when the facts are not clearly stated, they try to find the facts and approximate them, instead of analyzing what kind of information the public had at its disposal

\chapter{Narratives and Decision Making}


\section{What Narratives Are}

\subsection{McAdams Research on Narratives}

\subsection{Social Psychology Background}


\section{How Narratives can help}

\subsection{Bayesian Brain and Predictive Coding}

Here, there could be a direct link to the algorithms that are used in Machine Learning, AI, and NLP. 

\subsection{Influence and Change on Human Beings}

Akerlof and Shiller understand narratives as a convention, but it is more than that, it changes how people think and perceive the world. \cite{Akerlof.2016}


\section{Narrative Research}



\chapter{Natural Language Processing}

\section{Methods in Natural Language Processing}

\subsection{Sensitivity Analysis}

\subsection{Topic Modeling}


\chapter{Data and Methodology}

\chapter{Results}

\chapter{Conclusion}


\newpage

\phantomsection 
\addcontentsline{toc}{chapter}{Bibliography} 
\renewcommand\bibname{Bibliography}

\bibliography{EconomicNarrative}



\newpage

\appendix
\noappendicestocpagenum
\addappheadtotoc
%%\appendixpage



\renewcommand{\theequation}{A.\arabic{equation}}


\chapter{Whatever may come...}



\section{For Example... }

\renewcommand{\theequation}{B.\arabic{equation}}


\chapter{Whatever may come...}



\section{For Example... }





\newpage
{\pagestyle{firststyle}
	
\chapter*{Declaration of Authorship}

"I hereby declare
\begin{itemize}
	\item that I have written this thesis without any help from others and without the use of
	documents and aids other than those stated above;
	\item that I have mentioned all the sources used and that I have cited them correctly
	according to established academic citation rules;
	\item that I have acquired any immaterial rights to materials I may have used such as images
	or graphs, or that I have produced such materials myself;
	\item that the topic or parts of it are not already the object of any work or examination of
	another course unless this has been explicitly agreed on with the faculty member in
	advance and is referred to in the thesis;
	\item that I will not pass on copies of this work to third parties or publish them without the
	University’s written consent if a direct connection can be established with the
	University of St.Gallen or its faculty members;
	\item that I am aware that my work can be electronically checked for plagiarism and that I
	hereby grant the University of St.Gallen copyright in accordance with the Examination
	Regulations in so far as this is required for administrative action;
	\item that I am aware that the University will prosecute any infringement of this declaration
	of authorship and, in particular, the employment of a ghostwriter, and that any such
	infringement may result in disciplinary and criminal consequences which may result
	in my expulsion from the University or my being stripped of my degree."
\end{itemize}

%\noindent By submitting this academic term paper, I confirm through my conclusive action that I am
%submitting the Declaration of Authorship, that I have read and understood it, and that it is
%true.

%\hrule
\vspace*{3cm}
%\hrule


\dotbox{Location, Date} \hfill \dotbox{Signature}\\

\vspace*{.5 cm}
\noindent By submitting this academic term paper, I confirm through my conclusive action that I am
submitting the Declaration of Authorship, that I have read and understood it, and that it is
true.
\cleardoublepage
}
\end{document}
